\section*{Capítulo \ref{ch:g11:minerals} --- Los minerales}

\begin{solution}{exo:g11:minerals:1}
La \omterm{def:g11:minerals:cleavage}{exfoliación} es la tendencia de un mineral a romperse según planos lisos,
paralelos a direcciones en las que los enlaces entre sus átomos son más
débiles; la \omterm{def:g11:minerals:cleavage}{fractura} es la rotura según superficies irregulares, en un
mineral sin ningún plano de debilidad. El cuarzo no tiene \omterm{def:g11:minerals:cleavage}{exfoliación}:
muestra una \omterm{def:g11:minerals:cleavage}{fractura} concoidea, con superficies curvas y lisas como las del
vidrio roto. Una mica tiene una \omterm{def:g11:minerals:cleavage}{exfoliación}, perfecta: se separa en láminas.
\end{solution}

\begin{solution}{exo:g11:minerals:2}
Basta una traza de otro elemento para cambiar el color del cuarzo, que puede
ser incoloro, blanco, negro, rosa, azul o violeta; el olivino es casi siempre
verde. Más fiables: la \omterm{def:g11:minerals:cleavage}{exfoliación} (o la \omterm{def:g11:minerals:cleavage}{fractura}), la dureza, el \omterm{def:g11:minerals:lustre}{brillo}, el
\omterm{def:g11:minerals:habit}{hábito} de los cristales libres, la densidad.
\end{solution}

\begin{solution}{exo:g11:minerals:3}
(1) Olivino. (2) Biotita. (3) Ortosa (un feldespato: dos exfoliaciones,
dureza 6). (4) Cuarzo (sin \omterm{def:g11:minerals:cleavage}{exfoliación}, dureza 7).
\end{solution}

\begin{solution}{exo:g11:minerals:4}
(1) Cinco átomos: un silicio en el centro y cuatro oxígenos en los vértices de
un tetraedro. (2) $+4 + 4 \times (-2) = -4$: el grupo es \ce{SiO4^4-}.
\end{solution}

\begin{solution}{exo:g11:minerals:5}
Olivino: tetraedros aislados. Piroxenos: cadenas sencillas. Anfíboles:
cadenas dobles. Micas: láminas. Feldespatos y cuarzo: armazones.
\end{solution}

\begin{solution}{exo:g11:minerals:6}
(1) $M(\ce{SiO2}) = 28.085 + 2 \times 15.999 = \qty{60.083}{g/mol}$;
el oxígeno: $31.998/60.083 \approx \qty{53.3}{\percent}$.
(2) $M(\ce{Mg2SiO4}) = 2 \times 24.305 + 28.085 + 4 \times 15.999 =
\qty{140.691}{g/mol}$; el oxígeno: $63.996/140.691 \approx \qty{45.5}{\percent}$.
(3) La corteza, con un \qty{46}{\percent}, es del mismo orden: está formada
casi por completo por minerales ricos en oxígeno, los \omterm{def:g11:minerals:silicate}{silicatos}.
\end{solution}

\begin{solution}{exo:g11:minerals:7}
El grano blanco tiene dos exfoliaciones y una dureza entre 5.5 y 6.5: un
feldespato (6). El grano gris es más duro que 6.5 y muestra una \omterm{def:g11:minerals:cleavage}{fractura}
curva, concoidea, sin \omterm{def:g11:minerals:cleavage}{exfoliación}: cuarzo (7).
\end{solution}

\begin{solution}{exo:g11:minerals:8}
(1) $4 - 2/2 = 3$ oxígenos por silicio: \ce{SiO3}. (2) \ce{(Si,Al)2O6} tiene
dos átomos de silicio (o de aluminio) para seis oxígenos, tres por silicio: es
dos veces \ce{SiO3}.
\end{solution}

\begin{solution}{exo:g11:minerals:9}
$\rho = 0.55 \times 2.69 + 0.35 \times 3.4 + 0.10 \times 3.32 = 1.480 + 1.190
+ 0.332 \approx \qty{3.00}{g/cm^3}$. El cubo tiene un volumen de
\qty{1000}{cm^3}, luego una masa de unos \qty{3.00}{kg}.
\end{solution}

\begin{solution}{exo:g11:minerals:10}
(1) $0.75 \times 3.32 + 0.25 \times 3.4 = 2.49 + 0.85 = \qty{3.34}{g/cm^3}$.
(2) $(3.34 - 3.00)/3.00 \approx 0.11$: alrededor de un \qty{11}{\percent} más
densa. (3) El basalto y el \omterm{def:g11:ocean-ridges:gabbro}{gabro} de la \omterm{def:g11:ocean-ridges:oceanic-lithosphere}{corteza oceánica} son más ligeros que la
\omterm{def:g11:ocean-ridges:peridotite}{peridotita} del manto: se quedan encima de ella.
\end{solution}

\begin{solution}{exo:g11:minerals:11}
(1) Plagioclasa $168/400 = 0.42$, piroxeno $0.33$, olivino $0.15$, huecos
$0.10$. (2) La proposición vale para una roca sin poros; aquí una décima parte
del volumen está vacía. (3) En la parte sólida, las fracciones son $0.42/0.90
\approx 0.467$, $0.367$ y $0.167$:
$\rho_{\text{sólido}} \approx 0.467 \times 2.69 + 0.367 \times 3.4 + 0.167
\times 3.32 \approx \qty{3.06}{g/cm^3}$. Con sus huecos, que añaden volumen
pero no masa, la roca tiene la densidad $0.42 \times 2.69 + 0.33 \times 3.4 +
0.15 \times 3.32 \approx \qty{2.75}{g/cm^3}$, es decir, $0.9 \times 3.06$.
\end{solution}

\begin{solution}{exo:g11:minerals:12}
(1) Oxígeno $464/16.0 = 29.0$~mol; silicio $283/28.1 \approx 10.1$; aluminio
$84/27.0 \approx 3.1$; hierro $52/55.8 \approx 0.93$; calcio $46/40.1 \approx
1.15$; magnesio $28/24.3 \approx 1.15$; sodio $23/23.0 = 1.0$; potasio
$15/39.1 \approx 0.38$. (2) En total, unos $46.8$~mol de átomos, de los que
$29.0/46.8 \approx 0.62$: casi dos átomos de cada tres son de oxígeno.
(3) Los grandes iones óxido son la mayoría de los átomos, y ocupan una parte
todavía mayor del volumen: a la escala de los átomos, la corteza es un
apilamiento de iones óxido, con los pequeños iones de silicio y de los
metales en los huecos que quedan entre ellos.
\end{solution}

\begin{solution}{exo:g11:minerals:13}
(1) El número medio de vértices compartidos es $(2 + 3)/2 = 2.5$, así que hay
$4 - 2.5/2 = 2.75$ oxígenos por silicio. (2) Cuatro silicios para once
oxígenos: \ce{Si4O11}. (3) \ce{Si7AlO22} tiene ocho tetraedros (siete de
silicio, uno de aluminio) para 22 oxígenos, es decir, $22/8 = 2.75$ por
tetraedro, o dos veces \ce{Si4O11}. Los dos grupos \ce{OH} no forman parte de
los tetraedros.
\end{solution}

\begin{solution}{exo:g11:minerals:14}
$0.20 \times 2.62 + 0.10 \times 2.56 + 0.50 \times 2.65 + 0.10 \times 3.09 +
0.10 \times 3.23 = 0.524 + 0.256 + 1.325 + 0.309 + 0.323 \approx
\qty{2.74}{g/cm^3}$, por encima de \qty{2.7}{g/cm^3}; el granito del ejemplo,
más rico en cuarzo y en feldespato potásico, tiene unos \qty{2.59}{g/cm^3},
por debajo de \qty{2.6}{g/cm^3}. Las rocas del oeste son más densas porque
contienen más plagioclasa y más biotita y hornblenda, ricas en hierro y
magnesio; los granitos ricos en sílice contienen más cuarzo y feldespato
potásico, los minerales formadores de rocas más ligeros.
\end{solution}

\begin{solution}{exo:g11:minerals:15}
(1) Corteza superior: $66.6 \times 0.467 \approx \qty{31.1}{\percent}$ de
silicio; corteza entera: $60.6 \times 0.467 \approx \qty{28.3}{\percent}$.
(2) La corteza superior, más rica en sílice, debería contener más cuarzo y
feldespatos, y menos minerales ricos en hierro y magnesio; la corteza más
profunda, lo contrario. (3) El cuarzo y los feldespatos son los minerales
formadores de rocas más ligeros: la corteza superior debería ser más ligera
que la corteza más profunda.
\end{solution}

\begin{solution}{pb:g11:minerals:1}
\textbf{1.} Cuarzo: sin \omterm{def:g11:minerals:cleavage}{exfoliación}, \omterm{def:g11:minerals:cleavage}{fractura} concoidea, \omterm{def:g11:minerals:lustre}{brillo} vítreo, más
duro que el clavo (7).

\textbf{2.} Ortosa, un feldespato potásico: dos exfoliaciones, dureza 6 en la
escala de Mohs.

\textbf{3.} Plagioclasa (aquí, oligoclasa), el feldespato con sodio y
calcio.

\textbf{4.} D es biotita, y E, moscovita: dos micas. Sus tetraedros forman
láminas, y los enlaces entre las láminas son débiles: se separan a lo largo de
esa única \omterm{def:g11:minerals:cleavage}{exfoliación} perfecta.

\textbf{5.} El cuarzo y los dos feldespatos (A, B, C) son armazones; la
biotita (D) contiene hierro y magnesio.

\textbf{6.} Cuarzo $150/500 = 0.30$; ortosa $0.35$; oligoclasa $0.25$;
biotita $0.06$; moscovita $0.04$.

\textbf{7.} Cuando los granos están repartidos sin orden, la fracción de
puntos sobre un mineral estima la fracción de la superficie que cubre en la
cara, y esta es igual a la fracción del volumen que ocupa.

\textbf{8.} $0.35 + 0.25 = 0.60$: un \qty{60}{\percent} de feldespato.

\textbf{9.} En $1/500 = 0.002$, es decir, un \qty{0.2}{\percent}.

\textbf{10.} Sí: feldespato, cuarzo y mica.

\textbf{11.} $\rho = \phi_1 \rho_1 + \phi_2 \rho_2 + \dots + \phi_k \rho_k$,
con $\phi_1 + \dots + \phi_k = 1$.

\textbf{12.} Cuarzo $0.30 \times 2.62 = 0.786$; ortosa $0.35 \times 2.56 =
0.896$; oligoclasa $0.25 \times 2.65 = 0.6625$; biotita $0.06 \times 3.09
\approx 0.185$; moscovita $0.04 \times 2.82 \approx 0.113$ (todo en
\unit{g/cm^3}).

\textbf{13.} $0.786 + 0.896 + 0.6625 + 0.185 + 0.113 \approx
\qty{2.64}{g/cm^3}$.

\textbf{14.} La ortosa, con $0.896$, porque es la más abundante. La biotita y
la moscovita, con \qty{3.09}{} y \qty{2.82}{g/cm^3}, son más densas que los
demás minerales y elevan la densidad de la roca por encima de la del cuarzo y
el feldespato solos.

\textbf{15.} \qty{2640}{kg/m^3}: \qty{1}{m^3} del granito tiene una masa de
unas \qty{2.64}{t}.

\textbf{16.} $100 \times 2.643 \approx \qty{264.3}{g}$.

\textbf{17.} El bloque desaloja \qty{100}{cm^3} de agua, de masa
\qty{100}{g}: la balanza indica $264.3 - 100 = \qty{164.3}{g}$.

\textbf{18.} $264.3/(264.3 - 164.3) = 264.3/100 \approx \qty{2.64}{g/cm^3}$.

\textbf{19.} Entre los dos: más ligero que la media de \qty{2.69}{g/cm^3}, y
por encima de los \qty{2.6}{g/cm^3} de los granitos más ricos en sílice y
potasio. Más plagioclasa y más biotita u hornblenda, y menos cuarzo y
feldespato potásico, harían más denso un granito.

\textbf{20.} La densidad del granito es de unos
\textbf{\qty{2.64}{g/cm^3}}.
\end{solution}
